\PassOptionsToPackage{quiet}{xeCJK}
%% 2026 规范第十条:电子版 = electronic,第一页直接是摘要页
\documentclass[electronic,bwprint]{format}
\usepackage{etoolbox}
\usepackage{ctex}
\BeforeBeginEnvironment{tabular}{\zihao{-5}}
\usepackage[framemethod=TikZ]{mdframed}
\usepackage{url}
\usepackage{array}
\usepackage{tabularx}
\usepackage{longtable}
\newcolumntype{C}{>{\centering\arraybackslash}X}
\newcolumntype{R}{>{\raggedleft\arraybackslash}X}
\newcolumntype{L}{>{\raggedright\arraybackslash}X}

\renewcommand{\textbf}[1]{{\song\bfseries #1}}

\usepackage{tikz}
\usetikzlibrary{arrows.meta}
\tikzset{
  box/.style={rectangle, draw, minimum width=3.2cm, minimum height=0.9cm, align=center},
  arrow/.style={thick, -{Stealth}}
}

\title{NIPT 时点选择与胎儿异常判定的数学建模}

\begin{document}

%% electronic 模式:跳过承诺书和编号页,第一页直接是标题 + 摘要
\maketitle


\begin{abstract}

本文针对【题目】问题,基于【数据规模】的【数据类型】,综合运用【主要方法】,系统研究了【问题1】、【问题2】、【问题3】和【问题4】。

\textbf{针对问题一：}【简要描述问题1,采用方法,关键结果,关键数值】。

\textbf{针对问题二：}【简要描述问题2,采用方法,关键结果,关键数值】。

\textbf{针对问题三：}【简要描述问题3,采用方法,关键结果,关键数值】。

\textbf{针对问题四：}【简要描述问题4,采用方法,关键结果,关键数值】。

本文为【应用场景】提供了【主要贡献】的研究框架,研究结果对【实际意义】具有参考价值。

\keywords{关键词1;关键词2;关键词3;关键词4}
\label{abstract:end}

\end{abstract}

\thispagestyle{plain}

\newpage

\section{问题重述}

在撰写论文时,首先要简单地说明问题的情景,即要说清事情的来龙去脉。列出必要数据,提出要解决的问题,并给出研究对象的关键信息的内容,它的目的在于使读者对要解决的问题有一个印象。

% 该部分主要是用自己的话对问题进行总结,关键是改写!
% 一般问题重述的内容在半页到一页左右即可。
% 虽然该部分不是重点,但确实查重最容易出问题的部分!

\textbf{问题 1：}【问题1的简要描述,核心求解目标】。

\textbf{问题 2：}【问题2的简要描述,核心求解目标】。

\textbf{问题 3：}【问题3的简要描述,核心求解目标】。

\textbf{问题 4：}【问题4的简要描述,核心求解目标】。


\section{问题分析}

从实际问题到模型建立是一种从具体到抽象的思维过程,问题分析这一部分就是沟通这一过程的桥梁,因为它反映了建模者对于问题的认识程度如何,也体现了解决问题的雏形,起着承上启下的作用,也体现出建模者的综合水平。

% 这部分的内容应包括:题目中包含的信息和条件,利用信息和条件对题目做整体分析,确定用什么方法建立模型,一般是每个问题单独分析一小节,分析过程要简明扼要,不需要放结论。
% 建议在文字说明的同时用图形或图表(例如流程图)列出思维过程,这会使你的思维显得很清晰,让人觉得一目了然。

\subsection{问题一的分析}

【问题一是【问题类型:优化/评价/预测/分类/机理】问题,要求【具体求解目标】。从题目条件来看,【关键信息1】和【关键信息2】是核心约束,需要【建模思路概述】。本节将围绕【主要方法】展开,辅以【对比方法】进行交叉验证。】

\subsection{问题二的分析}

【问题二在问题一的基础上进一步考虑【新增约束/变量】,要求【具体求解目标】。需要【说明该问题的核心难点】,并采用【主要方法】解决。XXXXXXXXXXXXXXXXXXXX(一般写 10 行左右即可)】

\subsection{问题三的分析}

【问题三要求【具体求解目标】。XXXXXXXXXXXXXXXXXXXX(一般写 10 行左右即可)】

\subsection{问题四的分析}

【问题四是【建议/综合/扩展】型问题,基于前三问的结果,从【角度1】、【角度2】等维度给出【数量】条【建议/结论】。XXXXXXXXXXXXXXXXXXXX(一般写 10 行左右即可)】

针对以上问题,本文采用了以下整体思路框架进行系统性研究和分析:

\begin{figure}[ht]
  \begin{center}
    \begin{tikzpicture}[x=1cm, y=0.7cm]
      \node[box] (n11) at (0, 0) {问题重述};
      \node[box] (n12) at (5, 0) {问题分析};
      \node[box] (n13) at (10, 0) {模型假设};
      \node[box] (n21) at (0, -3) {数据预处理};
      \node[box] (n22) at (5, -3) {模型建立};
      \node[box] (n23) at (10, -3) {模型求解};
      \node[box] (n31) at (0, -6) {模型检验};
      \node[box] (n32) at (5, -6) {模型评价};
      \node[box] (n33) at (10, -6) {模型改进};
      \draw[arrow] (n11) -- (n12);
      \draw[arrow] (n12) -- (n13);
      \draw[arrow] (n13) -- ++(0,-1.0) -| (n21);
      \draw[arrow] (n21) -- (n22);
      \draw[arrow] (n22) -- (n23);
      \draw[arrow] (n23) -- ++(0,-1.0) -| (n31);
      \draw[arrow] (n31) -- (n32);
      \draw[arrow] (n32) -- (n33);
    \end{tikzpicture}
    \caption{问题的总分析}
    \label{fig:overall_analysis}
  \end{center}
\end{figure}


\section{模型假设}

针对上述研究问题,本文在建模过程中作出以下合理假设:

\textbf{假设 1：}【假设1的描述,关于数据来源、样本代表性等】。

\textbf{假设 2：}【假设2的描述,关于变量关系的连续性、近似条件等】。

\textbf{假设 3：}【假设3的描述,关于数据独立性、个体内相关性等】。

\textbf{假设 4：}【假设4的描述,关于测量误差、噪声分布等】。

\textbf{假设 5：}【假设5的描述,关于风险函数、目标函数的假设】。

\textbf{假设 6：}【假设6的描述,关于模型推广性、适用范围的说明】。


\section{符号说明}

\begin{longtable}{>{\centering\arraybackslash}p{0.20\textwidth}>{\centering\arraybackslash}p{0.20\textwidth}>{\centering\arraybackslash}p{0.54\textwidth}}
  \caption{本文主要符号说明}\\
  \toprule
  符号 & 含义 & 说明 \\
  \midrule
  \endfirsthead
  \caption{本文主要符号说明(续)}\\
  \toprule
  符号 & 含义 & 说明 \\
  \midrule
  \endhead
  \bottomrule
  \endfoot
  $Y$ & 【响应变量】 & 【说明】\\
  $T$ & 【自变量1】 & 【说明】\\
  $X$ & 自变量向量 & 【说明】\\
  $r$ & Pearson 相关系数 & 【说明】\\
  $\beta_j$ & 回归系数 & 【说明】\\
  $p_j$ & 显著性 $p$ 值 & 【说明】\\
  $R^2$ & 决定系数 & 【说明】\\
  $F$ & F 统计量 & 【说明】\\
  $p$ & 概率 & 【说明】\\
  $R$ & 风险函数 & 【说明】\\
  $\alpha, \beta, \gamma$ & 【参数】 & 【说明】\\
  $Z$ & Z 值 & 【说明】\\
  $C$ & 【指标】 & 【说明】\\
  $N$ & 数量 & 【说明】\\
  $\mathrm{OR}$ & 优势比 & 【说明】\\
\end{longtable}


%%%%%%%%%%%%%%%%%%%%%%%%%%%%%%%%%%%%%%%%%%%%%%%%%%%%%%%%%%%%%
\section{模型的建立与求解}


\subsection{问题一的模型的建立和求解}

\subsubsection{具体分析}

问题X是【问题类型:优化/评价/预测/分类/机理】问题,要求【具体求解目标】。

本问采用【主要方法】作为主体方法,辅以【对比方法】进行交叉验证。具体地,首先对【数据/对象】进行预处理,然后【建模步骤1】,接着【建模步骤2】,最后【求解与验证】。

具体分析流程如图\ref{fig:analysis1_flow}所示。

\begin{figure}[ht]
  \begin{center}
    \begin{tikzpicture}[x=1cm, y=0.7cm]
      \node[box] (n11) at (0, 0) {数据预处理};
      \node[box] (n12) at (5, 0) {【步骤1】};
      \node[box] (n13) at (10, 0) {【步骤2】};
      \node[box] (n22) at (5, -3) {【对比方法】};
      \node[box] (n21) at (0, -3) {主体方法};
      \node[box] (n23) at (10, -3) {【验证】};
      \node[box] (n31) at (0, -6) {结果分析};
      \node[box] (n32) at (5, -6) {灵敏度};
      \node[box] (n33) at (10, -6) {结论};
      \draw[arrow] (n11) -- (n12);
      \draw[arrow] (n12) -- (n13);
      \draw[arrow] (n13) -- ++(0,-1.0) -| (n23);
      \draw[arrow] (n23) -- (n22);
      \draw[arrow] (n22) -- (n21);
      \draw[arrow] (n21) -- (n31);
      \draw[arrow] (n31) -- (n32);
      \draw[arrow] (n32) -- (n33);
    \end{tikzpicture}
    \caption{问题X分析流程图}
    \label{fig:analysis1_flow}
  \end{center}
\end{figure}

\subsubsection{模型准备}

在数据预处理的基础上,选取【响应变量】,自变量包括【因素1、因素2、...】。对所有数值变量进行标准化处理:
\begin{equation}
x'_j = \frac{x_j - \bar{x}_j}{s_j}
\end{equation}
其中,$\bar{x}_j$ 和 $s_j$ 分别为第 $j$ 个变量的样本均值和标准差。

该问题本质上是一个【问题类型】问题。常用的【方法类别】包括【方法A】、【方法B】和【方法C】。本节从【维度1】、【维度2】和【维度3】三方面进行评估。

\begin{longtable}{>{\centering\arraybackslash}p{0.20\textwidth}>{\centering\arraybackslash}p{0.20\textwidth}>{\centering\arraybackslash}p{0.20\textwidth}>{\centering\arraybackslash}p{0.18\textwidth}>{\centering\arraybackslash}p{0.18\textwidth}}
  \caption{【方法类别】对比}\\
  \toprule
  方法 & 【维度1】 & 【维度2】 & 【维度3】 & 【维度4】 \\
  \midrule
  \endfirsthead
  \caption{【方法类别】对比(续)}\\
  \toprule
  方法 & 【维度1】 & 【维度2】 & 【维度3】 & 【维度4】 \\
  \midrule
  \endhead
  \bottomrule
  \endfoot
  方法A & 【评级】 & 【评级】 & 【评级】 & 【评级】 \\
  方法B & 【评级】 & 【评级】 & 【评级】 & 【评级】 \\
  方法C & 【评级】 & 【评级】 & 【评级】 & 【评级】 \\
\end{longtable}

通过对比可以看出,【方法A】在【某维度】方面表现优异,【方法C】在【某维度】独树一帜。综合考量,本文采用【选定方法组合】的策略,【说明理由】。

综上所述,本节完成了数据预处理和分析方法选型,接下来将进行【下一步内容】。


\subsubsection{模型建立}

设 $Y_i$ 为第 $i$ 个【观测对象】的【响应变量】,自变量向量 $X_i = (X_{i1}, X_{i2}, \ldots, X_{ip})^T$ 分别对应【因素说明】。建立【模型类型】:
\begin{equation}
Y_i = f(X_i; \theta) + \varepsilon_i
\end{equation}
其中,$\theta$ 为待估参数向量,$\varepsilon_i$ 为随机误差项。

为便于比较各因素的影响程度,对自变量进行标准化处理后拟合模型。对各参数采用最大似然估计或最小二乘法,通过适当的统计检验判断其显著性。

最后,采用【对比方法】作为补充,提供【特征重要性 / 非线性关系】等信息。

\subsubsection{模型求解}

基于【样本数】个【样本说明】,采用上述方法进行求解,关键结果如下。

【简要描述主要发现】。单变量分析表明,【因素1】与响应变量呈显著【正/负】相关($r = $ 【值】, $p = $ 【值】)。

【主方法】结果显示,【关键指标】达到【值】,【次要指标】为【值】。各因素对响应变量的影响系数如图\ref{fig:q1_coef}所示。

\begin{figure}[ht]
  \begin{center}
    \includegraphics[width=0.85\textwidth]{../求解/问题1/图片/4.回归系数.png}
    \caption{【主方法】系数}
    \label{fig:q1_coef}
  \end{center}
\end{figure}

【对比方法】的特征重要性排序为:【因素1】 【值】,【因素2】 【值】,【因素3】 【值】。前两个因素在两种方法中均排名前二,验证了它们对【响应变量】的影响,如图\ref{fig:q1_rfimp}所示。

\begin{figure}[ht]
  \begin{center}
    \includegraphics[width=0.85\textwidth]{../求解/问题1/图片/5.随机森林特征重要性.png}
    \caption{【对比方法】特征重要性}
    \label{fig:q1_rfimp}
  \end{center}
\end{figure}

5 折交叉验证结果显示,【指标】为【值 ± 标准差】,【进一步说明】。

预测值与实际值的散点图如图\ref{fig:q1_pred}所示。

\begin{figure}[ht]
  \begin{center}
    \includegraphics[width=0.85\textwidth]{../求解/问题1/图片/6.预测vs实际.png}
    \caption{预测值与实际值散点图}
    \label{fig:q1_pred}
  \end{center}
\end{figure}

\subsubsection{结果分析}

研究结果表明,【关键结论】。这一发现与【领域直觉/理论】一致,【进一步解释】。

【进一步说明】。然而,【简单方法】对【响应变量】变异的解释力有限(【指标】较低),提示实际关系中存在复杂的非线性成分和未观测的混杂因素。

综上所述,本问建立了【响应变量】与【自变量】的【模型类型】关系模型,识别出【主要因素】是主要影响因素,模型【整体显著/解释力有限】,需要更复杂的建模方法以提高拟合优度。



\subsection{问题二的模型的建立和求解}
% 问题 2 模板 - 复制 5.1.1 内容并修改
% 用户使用时请将 5.1.1.分析与准备.tex 内容复制到此,然后修改
\input{5.1.1.分析与准备.tex}

% 问题 2 模板 - 复制 5.1.2 内容并修改
% 用户使用时请将 5.1.2.建模与求解.tex 内容复制到此,然后修改
\input{5.1.2.建模与求解.tex}



\subsection{问题三的模型的建立和求解}
% 问题 3 模板 - 复制 5.1.1 内容并修改
\input{5.1.1.分析与准备.tex}

% 问题 3 模板 - 复制 5.1.2 内容并修改
\input{5.1.2.建模与求解.tex}



\subsection{问题四的模型的建立和求解}
\input{5.4.1.分析与准备.tex}
\input{5.4.2.建模与求解.tex}



\section{模型检验}

本节对前文建立的主要模型进行检验,包括误差分析、交叉验证和灵敏度分析三个方面。

\subsection{问题一模型的检验}

对于问题一的【主方法】,采用【5 折 / 10 折】交叉验证评估其泛化能力。结果显示,【指标】为【值 ± 标准差】,【进一步说明】。

\subsection{问题二模型的检验}

对问题二的【主方法】,采用交叉验证和【误差分析方法】两种方式进行检验。

【描述检验方法与结果】。

\subsection{问题三模型的检验}

对问题三的【主方法】,从以下【N】个角度进行检验。

\textbf{交叉验证:}【结果说明】。

\textbf{【其他检验方法】:}【结果说明】。

\textbf{【灵敏度分析】:}对【关键因素】进行 $\pm 20\%$ 缩放,观察【响应指标】变化。结果显示,【最敏感因素】对【结果】影响最显著。

\subsection{问题四模型的检验}

对问题四的【分类/预测】模型,主要采用 5 折分层交叉验证和【其他评估方法】进行检验。

\textbf{5 折交叉验证:}【主方法】的交叉验证准确率【值】, F1 分数【值】;【对比方法】的交叉验证准确率【值】, F1 分数【值】。

\textbf{【其他方法】:}【结果说明】。

\textbf{类别不平衡处理:}【说明处理方法及效果】。

综合上述检验,本文所建立的主要模型在交叉验证和误差传播分析中均表现出较好的稳健性,【进一步说明】。


\section{模型优缺点评价}

% 所谓的模型优缺点评价往往并不局限于模型本身,在整个建模过程中所表露出的优缺点均可在最后进行陈述,一般撰写模型优缺点的基本原则是优点说充分,缺点不回避。

\subsection{模型的优点}

常见的优点表述形式如下:

\begin{itemize}[itemindent=2em]
  \item \textbf{优点 1：模型或思路设计的简洁实用,效率高。}【结合本文说明】~\cite{ref1}。
  \item \textbf{优点 2：本文建立的模型具有很强的创新性。}【结合本文说明】。
  \item \textbf{优点 3：模型的计算结果准确,精度高。}【结合本文说明】。
  \item \textbf{优点 4：模型考虑的系统全面,有很强的实用性。}【结合本文说明】。
  \item \textbf{优点 5：对模型进行了各类检验、稳定性高。}【结合本文说明】。
  \item \textbf{优点 6：模型本身具有的优点。}【结合本文说明】。
\end{itemize}

\subsection{模型的缺点}

常见的缺点表述形式如下:

\begin{itemize}[itemindent=2em]
  \item \textbf{缺点 1：}受【XX 因素】限制,未考虑【XX 情况】,影响精度。
  \item \textbf{缺点 2：}本文考虑的因素较为理想,降低了模型的普适性和推广能力。
  \item \textbf{缺点 3：}由于系统考虑了【XXX】等因素,导致模型较为复杂,计算时间长,效率低。
  \item \textbf{缺点 4：}模型本身具有的缺点。
\end{itemize}

\subsection{模型的改进}

% 模型的改进一般是针对模型的缺点而言的,主要是提出一些改进的思路即可。

针对上述不足,本文可从以下方面进行改进:【改进方向 1】,通过【具体方法】以解决【具体问题】;【改进方向 2】,通过【具体方法】以提升【具体指标】;【改进方向 3】,通过【具体方法】以扩展【应用场景】。后续研究将围绕上述方向进一步深入,使模型在更广泛的实际场景中保持高精度与强鲁棒性。


%% 2026 AI 规定第 3 条: AI 工具使用声明放在参考文献之前
\newpage

% AI 工具使用声明
% 依据《全国大学生数学建模竞赛人工智能工具使用规定》(2026 试行) 第 3 条
% 必须放在 9.参考文献.tex 之前
\section*{AI 工具使用声明}

\begin{quote}
本参赛队在竞赛过程中使用了 AI 工具,主要用于 \textbf{【简要用途:语言润色、代码调试、模型分析、论文排版】},详细使用情况见支撑材料中的 ``AI 工具使用详情.pdf''。
\end{quote}

\vspace{0.5em}

\noindent\textbf{具体使用情况概述:}

\begin{itemize}
  \item \textbf{AI 工具名称:} 【请填写所用 AI 工具名称,如:大语言模型 X】
  \item \textbf{版本或型号:} 【请填写版本号】
  \item \textbf{使用目的:} 【简要描述,如辅助完成问题分析、Python 求解代码编写、LaTeX 论文排版等】
  \item \textbf{使用环节:} 问题分析 $\rightarrow$ 模型比选 $\rightarrow$ Python 求解 $\rightarrow$ LaTeX 论文写作 $\rightarrow$ 验收修改
  \item \textbf{人工审查与核验:} 全部 AI 生成的代码均在本地 Python 环境下运行验证,生成的数学模型、公式推导和结论均经人工复核;AI 生成的 LaTeX 文本经过逐字人工校对与修改,所有数据结果与原始代码运行结果一致
\end{itemize}

\noindent\textbf{关于核心建模的声明:} 本论文的核心建模思路、模型选择、参数设置、结果分析均由参赛队主导完成,AI 工具仅作为辅助手段,所有最终结论由参赛队独立负责。

\noindent\textbf{关于参赛队责任的声明:} 依据《全国大学生数学建模竞赛参赛规则》(2026 修订稿) 第 6 条,本参赛队对所提交作品的**原创性、真实性和准确性负全部责任**。若论文中存在任何由 AI 工具辅助生成但未充分核验的内容,或与 AI 实际参与情况不符的描述,均由本参赛队承担相应责任。

\vspace{1em}

\noindent 详细使用情况记录在支撑材料中的 ``AI 工具使用详情.pdf'' 文件中,内容应包括:
\begin{enumerate}
  \item 所用 AI 工具名称、版本或型号
  \item 具体使用目的和环节
  \item 主要提示方式与使用过程说明(含典型交互示例)
  \item 对 AI 输出的采纳、人工修改和核验情况
\end{enumerate}


\newpage

% 参考文献
% 优先引用中文文献,8-15 条,GB/T 7714-2015 格式,每条在正文有引用。
%
% 对于常见的各类参考文献标注方法如下:
%   1) 著作:作者姓名,题名[M].出版地:出版社,出版年.
%   2) 期刊论文:作者姓名.题名[J].期刊名称,年,卷(期):页码.
%   3) 会议论文集:作者姓名.题名[C]//论文集名称,会议地点,会议日期.
%   4) 学位论文:作者姓名.题名[D].出版地:出版者,出版年.
%   5) 专利文献:专利申请者或所有者姓名.专利题名:专利国别,专利号[P].公告日期或公开日期.获取路径.
%   6) 电子文献:作者姓名.题名[文献类型标志(含文献载体标志)见其它].出版地:出版者,出版年(更新或修改日期),获取路径.
%   7) 报告:作者姓名.题名[R].出版地:出版者,出版年.
%   8) 标准:标准号.题名[S].出版地:出版者,出版年.
%
% 注意使用 AI 工具生成的内容,应在正文相应位置进行标注;
% 在参考文献中列出所用 AI 工具,格式如下:
%   [编号]工具名称,版本/型号,开发机构/公司,使用日期。
%   示例:[1]DeepSeek,DeepSeek-R1-0528,深度求索(DeepSeek),2025-09-05

\begin{thebibliography}{99}

% ========== 中文文献(优先) ==========
\bibitem{ref1} 【作者1】, 【作者2】. 【文献标题】[J]. 【期刊名】, 【年】, 【卷(期)】: 【起止页】.

\bibitem{ref2} 【作者1】, 【作者2】. 【文献标题】[M]. 【出版社所在地】: 【出版社】, 【年】.

\bibitem{ref3} 【作者1】, 【作者2】. 【文献标题】[D]. 【所在地】: 【授予单位】, 【年】.

\bibitem{ref4} 【作者1】, 【作者2】. 【文献标题】[C]//【会议名】. 【会议地点】: 【出版社】, 【年】: 【起止页】.

\bibitem{ref5} 【作者1】, 【作者2】. 【文献标题】[R]. 【所在地】: 【机构】, 【年】.

\bibitem{ref6} 【作者1】, 【作者2】. 【文献标题】[J]. 【期刊名】, 【年】, 【卷(期)】: 【起止页】.

\bibitem{ref7} 【作者1】, 【作者2】. 【文献标题】[J]. 【期刊名】, 【年】, 【卷(期)】: 【起止页】.

\bibitem{ref8} 【作者1】, 【作者2】. 【文献标题】[M]. 【出版社所在地】: 【出版社】, 【年】.

\bibitem{ref9} 【作者1】, 【作者2】. 【文献标题】[D]. 【所在地】: 【授予单位】, 【年】.

\bibitem{ref10} 【作者1】, 【作者2】. 【文献标题】[J]. 【期刊名】, 【年】, 【卷(期)】: 【起止页】.

% ========== AI 工具(必列) ==========
% 格式:[编号]工具名称,版本/型号,开发机构/公司,使用日期。
\bibitem{aideepseek} DeepSeek, DeepSeek-R1-0528, 深度求索(DeepSeek), 2025-09-05.

\end{thebibliography}


\newpage

\section*{附录}

\subsubsection*{一、程序使用说明}

\textbf{本论文使用了程序。}建模与求解所用到的全部源程序代码、关键中间图与结果数据详见下表。

\begin{longtable}{>{\centering\arraybackslash}p{0.40\textwidth}>{\centering\arraybackslash}p{0.56\textwidth}}
  \caption{附件说明表}\\
  \toprule
  文件 & 说明 \\
  \midrule
  \endfirsthead
  \caption{附件说明表(续)}\\
  \toprule
  文件 & 说明 \\
  \midrule
  \endhead
  \bottomrule
  \endfoot
  求解/预处理数据.py & 通用预处理脚本 \\
  求解/预处理数据.csv & 预处理后共享数据 \\
  求解/问题1/问题1\_xxx.py & 问题 1 求解代码 \\
  求解/问题1/图片/ & 问题 1 求解过程图 \\
  求解/问题1/结果/ & 问题 1 求解结果 csv \\
  求解/问题2/问题2\_xxx.py & 问题 2 求解代码 \\
  求解/问题2/图片/ & 问题 2 求解过程图 \\
  求解/问题2/结果/ & 问题 2 求解结果 csv \\
  求解/问题3/问题3\_xxx.py & 问题 3 求解代码 \\
  求解/问题3/图片/ & 问题 3 求解过程图 \\
  求解/问题3/结果/ & 问题 3 求解结果 csv \\
  求解/问题4/问题4\_xxx.py & 问题 4 求解代码 \\
  求解/问题4/图片/ & 问题 4 求解过程图 \\
  求解/问题4/结果/ & 问题 4 求解结果 csv \\
\end{longtable}

\noindent\textbf{运行说明:} 全部 Python 程序均基于 Python 3.12 编写,主要依赖 numpy、pandas、scikit-learn、scipy、matplotlib 等开源库。在已配置好依赖的环境中,直接运行各问题对应的 \texttt{问题X\_xxx.py} 即可复现论文中的全部图表与结果。预处理脚本 \texttt{预处理数据.py} 需在读取原始数据后首先执行,其输出 \texttt{预处理数据.csv} 供后续各问题共享。

\subsubsection*{二、支撑材料说明}

\textbf{本论文提供了支撑材料。}支撑材料内容包括:原始数据附件、问题求解完整源程序、问题求解结果 CSV、中间过程图(PNG)与 AI 工具使用详情(PDF)。具体清单如下表所示。

\begin{longtable}{>{\centering\arraybackslash}p{0.40\textwidth}>{\centering\arraybackslash}p{0.56\textwidth}}
  \caption{支撑材料文件清单}\\
  \toprule
  文件 & 说明 \\
  \midrule
  \endfirsthead
  \caption{支撑材料文件清单(续)}\\
  \toprule
  文件 & 说明 \\
  \midrule
  \endhead
  \bottomrule
  \endfoot
  数据/附件.xlsx & 原始 NIPT 检测数据 \\
  求解/预处理数据.py & 通用预处理脚本 \\
  求解/预处理数据.csv & 预处理后共享数据 \\
  求解/求解计划.md & 方案比选与详细求解计划 \\
  求解/问题*/图片/*.png & 全部中间过程图 \\
  求解/问题*/结果/*.csv & 全部结果 CSV \\
  求解/问题*/问题*\_*.py & 全部源程序 \\
  论文/AI工具使用详情.pdf & AI 工具使用详情 \\
\end{longtable}

\noindent\textbf{提交方式:} 全部支撑材料文件应使用 WinRAR 或类似软件压缩为一个文件(后缀为 RAR 或 ZIP,大小不超过 20MB),与论文电子版(PDF)分开提交。承诺书和编号专用页不要放在支撑材料中。所有文件中不能有显示参赛者身份和所在学校及赛区的信息。

\subsubsection*{三、运行环境}

本文的建模和可视化工作基于以下环境完成:

\begin{itemize}
  \item Python 3.12
  \item numpy, pandas, scikit-learn, scipy, matplotlib 等开源库
  \item LaTeX 编译环境(支持 ctex 宏包)
  \item xelatex 编译器(必需,format.cls 依赖)
\end{itemize}


\end{document}
